\chapter*{Glossaire}
\addcontentsline{toc}{chapter}{Glossaire}

\begin{description}
  \item[FCFA] Franc de la Communauté Financière Africaine, monnaie utilisée au Sénégal.
  \item[Modèle hédonique] Modèle économétrique postulant que le prix d'un bien composite (ex. un logement) se décompose en la somme des valeurs implicites de ses caractéristiques (surface, localisation, etc.).
  \item[RMSE] \emph{Root Mean Squared Error} : racine de l'erreur quadratique moyenne, métrique d'évaluation pénalisant fortement les grandes erreurs.
  \item[MAE] \emph{Mean Absolute Error} : erreur absolue moyenne entre les valeurs prédites et réelles.
  \item[MAPE] \emph{Mean Absolute Percentage Error} : erreur moyenne exprimée en pourcentage de la valeur réelle.
  \item[$R^2$] Coefficient de détermination : part de la variance de la variable cible expliquée par le modèle.
  \item[Random Forest] Algorithme d'apprentissage ensembliste combinant un grand nombre d'arbres de décision entraînés sur des sous-échantillons aléatoires des données.
  \item[XGBoost] \emph{Extreme Gradient Boosting} : algorithme de boosting de gradient optimisé, construisant les arbres de façon séquentielle pour corriger les erreurs des précédents.
  \item[MLP] \emph{Multi-Layer Perceptron} : réseau de neurones artificiels à propagation avant, composé d'une ou plusieurs couches cachées.
  \item[Hyperparamètre] Paramètre de configuration d'un modèle fixé avant l'entraînement (ex. profondeur d'un arbre, taux d'apprentissage), par opposition aux paramètres appris à partir des données.
  \item[One-Hot Encoding] Technique d'encodage transformant une variable catégorielle en un ensemble de variables binaires indicatrices.
  \item[Web scraping] Technique de collecte automatisée de données publiées sur des pages web.
\end{description}
